\documentclass[11pt]{article}
\usepackage[margin=1.1in]{geometry}
\usepackage{amsmath,amssymb,amsthm}
\usepackage{hyperref}
\newtheorem{proposition}{Proposition}
\newtheorem{remark}[proposition]{Remark}
\renewcommand{\Re}{\operatorname{Re}}
\title{A Lean-checked, counts-independent form of tomoto0's floor-class barrier in the cubic-theta zero-free half-plane argument}
\author{S.~Reiner\thanks{The exponent bookkeeping was carried out with the assistance of Claude (Anthropic), by two independent model instances reading the manuscripts, and cross-checked against Liu's formalized constants. The author takes responsibility for the content.}}
\date{October 28, 2026 --- draft v6 (priority corrected: the barrier is tomoto0's, 2026-10-07), not yet externally reviewed}
\begin{document}
\maketitle

\begin{abstract}
OpenAI's cubic-theta argument gives a uniform zero-free half-plane $\Re s>7/8$ for all finite-order Hecke $L$-functions over $\mathbb Q(\sqrt{-3})$ and all Dirichlet $L$-functions; Liu optimised its free parameters and proved that $B_{\mathrm{new}}=0.874957\ldots$ is the infimum of the resulting exponent comparison system. The barrier set by the \emph{floor class} was identified by tomoto0 in a barrier analysis of 2026-10-07 \cite{Tom26}: with the paper's normalisation $M+\ell=1$, the trivially counted rows near the critical line force $\sigma_0>13/15$ even with density-hypothesis-strength counts, and ``including the detector floor $a=51/100$, the floor alone gives $\sigma_0\ge0.86979$''. The present note records a sharpened and machine-checked form of that observation: within Liu's comparison system (free $M$, arbitrary witness counts), the bound is exactly $(13+18\delta_0)/(3(5+6\delta_0))=167/192$, obtained as a convex combination of three affine constraints, and formalized in Lean~4. The floor class is Poisson rows for which the class definition selects no zero witness (essentially, no zero with real part above the detector floor $a_0=51/100$), and which are therefore counted trivially. We also show that the obstruction cannot be removed by any estimate passing through absolute values of the scalar quotient of $L$-values, and state the signed cancellation that would be required. Nothing here is a priority claim over \cite{Tom26}; our literature search missed that report, which is written in Japanese and hosted on GitHub.
\end{abstract}

\section{Prior work: tomoto0's barrier analysis}
[LIT] tomoto0 \cite{Tom26} (2026-10-07, in Japanese), in the repository that also contains the independent Lean replay of the $7/8$ theorem, derives $\sigma_0=1-h/6+b/12$ from the manuscript's (15.9), notes that under $M+\ell=1$ this is $2/3+l_x/2+b/4$, and shows (``Barrier B'') that rows with zeros near the critical line, which carry the trivial count $R(0)=1$, force $\sigma_0>13/15$ even if all other counts were as strong as the density hypothesis; the detector floor $a=51/100$ alone gives $\sigma_0\ge0.86979$. The same report finds the optimum $\approx0.87496$ of the paper's bookkeeping and the endpoint margin $2.28\cdot10^{-4}$. The tracker \cite{Tracker} records $13/15$ as ``Limit of OpenAI's method'' on this basis. Everything below should be read as a formal restatement of Barrier B inside Liu's system, with the exact rational value and a proof checker.

\section{Liu's comparison system}

We use the notation of Liu \cite{Liu26}, who follows the 30 September 2026 manuscript \cite{Ope26}. The argument compares a probe $J_\eta(Z)$ with a Mellin signal containing $1/L^S(s,\eta)$; a direct bound $|J_\eta|\ll Z^{C(\sigma_0)+\omega}$ and a comparison bound $|J_\eta-f_\eta|\ll Z^{C(\beta_*)-\sigma}$ contradict $\beta_*>\sigma_0$ \cite[Prop.~2.1]{Ope26}. The free length parameters are $b=l_y-l_x\ge0$, the slot length $\ell$, and $M=l_x+l_y$, with $h:=1-l_x+\ell$. The direct bound has exponent $l_x/2+b/12$ and the principal Mellin exponent is $C_b(s)=s-(4+b)/6$; with $M+\ell=1$ this gives the boundary $B=11/12-\ell/4$ \cite[(30.2)]{Liu26}, and $\ell=1/6$ is the $7/8$ of \cite{Ope26}.

Liu's Proposition~25.1 isolates three affine constraints that force $B\ge B_{\mathrm{alg}}=(1507-2\sqrt{921})/1653$:
\begin{equation}\label{eq:L}
B\ge L_0:=\tfrac56+\tfrac M{12}-\tfrac\ell6,\qquad B\ge L_K:=\tfrac13+\tfrac{7M}{12}+\tfrac\ell3,
\end{equation}
from the two branches of the direct (reflected-energy) estimate, and
\begin{equation}\label{eq:H}
B\ge H_c:=1+\delta_c-(1+\delta_c)\tfrac M2+\delta_c\ell,\qquad \delta_c=\tfrac{49-\sqrt{921}}{48},
\end{equation}
from the critical zero class, where the count exponent is $R^*=2/3$. The general nonprincipal-row constraint behind \eqref{eq:H} is, from \cite[(18.4)]{Liu26} with $a=(1+\delta)/2$, $d=h$, $l_y=M-l_x$, $l_x=1+\ell-h$ (his (34.8) is the case $x=1/2$),
\begin{equation}\label{eq:E}
E_B(h;\delta,x)=H(\delta,x)-B+h\bigl(R(\delta,x)-\tfrac23\bigr)\le0,\qquad H(\delta,x)=1+\delta-(1+\delta)\tfrac M2+\delta\ell\bigl(\tfrac12+x\bigr),
\end{equation}
for every zero class $\delta=2a-1\in[\delta_0,5/6]$ and amplitude $x\in[0,1/2]$, with $R(\delta,x)$ the proven row-count exponent and $\delta_0=1/50$ the detector floor ($a_0=51/100$). Since $b=M-2(1+\ell-h)$, the constraint $b\ge0$ reads $h\ge1-M/2+\ell$; Liu's admissible set has the strict forms $b>0$ and $l_x>\ell$ (hence $h<1$), of which we use the closures.

\section{The floor class}

Rows in the lowest zero class $a=a_0$ are bounded with the trivial count $R=1$ and amplitude $\bar g\le\delta_0/2$, i.e.\ $x\le1/2$ \cite[Lemma~19.2(4)]{Liu26}; since \eqref{eq:E} is increasing in $x$, the binding case is $x=1/2$ (cf.\ the display $E_{B,\mathrm{floor}}$ in the proof of Lemma~34.3 of \cite{Liu26}), giving
\begin{equation}\label{eq:F}
B\ \ge\ F:=1+\delta_0-(1+\delta_0)\tfrac M2+\delta_0\ell+\tfrac h3 .
\end{equation}
This constraint does not involve any zero-witness count, so it is unaffected by any improvement of the estimates that produce $R(\delta,x)$ for $\delta>\delta_0$. Liu already notes, in the proof of his Lemma~34.3, that the trivial-count comparison for a raised floor requires $B\ge H_{\xi}+h/3$; what is added here is only the convex combination with the direct constraints, which makes the resulting bound independent of every count and of $(M,\ell,h)$. What is held fixed below is exactly: (i) the direct majorant and principal exponent of \cite[(22.1), (22.4)]{Liu26}, which give \eqref{eq:L}; (ii) the row exponent \cite[(18.4)]{Liu26} with the trivial count for the class $a_0$, which gives \eqref{eq:F}; (iii) $b\ge0$. Sharpening any of these (for instance the reflected-energy estimate behind (22.4), or the bound for the $a_0$ class by a means other than its count) changes the system.

\begin{proposition}\label{prop:barrier}
Let $M,\ell,h,B$ be real numbers satisfying \eqref{eq:L}, \eqref{eq:F} and $h\ge1-M/2+\ell$. Then
\[
B\ \ge\ \frac{13+18\delta_0}{3(5+6\delta_0)} ,
\]
which equals $167/192$ for $\delta_0=1/50$ and tends to $13/15$ as $\delta_0\to0$. Equality holds at $h=1-M/2+\ell$ (i.e.\ $b=0$), $\ell=(1-2\delta_0)/(5+6\delta_0)$, $M=1-\ell$ ; for $\delta_0=1/50$ this is $M=13/16$, $\ell=3/16$, $h=25/32$.
\end{proposition}
\begin{proof}
Since $F$ is increasing in $h$, \eqref{eq:F} with $h\ge1-M/2+\ell$ gives $B\ge F_{\min}:=1+\delta_0-(1+\delta_0)M/2+\delta_0\ell+(1-M/2+\ell)/3$. The positive weights
\[
w_0=\frac{2(5+9\delta_0)}{3(5+6\delta_0)},\qquad w_K=\frac{2}{3(5+6\delta_0)},\qquad w_F=\frac1{5+6\delta_0}
\]
sum to one, and in $w_0L_0+w_KL_K+w_FF_{\min}$ the coefficients of $M$ and of $\ell$ vanish identically; the constant that remains is $(13+18\delta_0)/(3(5+6\delta_0))$. Hence $B\ge\max\{L_0,L_K,F_{\min}\}\ge w_0L_0+w_KL_K+w_FF_{\min}$ gives the bound. For $\delta_0=1/50$ the weights are $259/384$, $25/192$, $25/128$ and the constant is $167/192$. Equality is checked by substitution. The inequality and its equality case are formalized in Lean~4 with Mathlib (file \texttt{FloorBarrier.lean} accompanying this note; theorems \texttt{barrier\_general} and \texttt{barrier\_167\_192}, depending only on the axioms \texttt{propext}, \texttt{Classical.choice}, \texttt{Quot.sound}). At the equality point $M+\ell=1$ and $b=0$, so the barrier is simply $L_0=11/12-\ell/4$ against $F=2/3+\delta_0/2+(1+3\delta_0/2)\ell$ on Liu's balanced line; $L_K$ only enforces $M+\ell\le1$.
\end{proof}

\begin{remark}[what the witness counts can and cannot do]
Replacing $R(\delta,x)$ by $R(\delta,x)-\rho$ uniformly in \eqref{eq:E} and minimising $B$ over $(M,\ell,h)$ subject to \eqref{eq:L}, \eqref{eq:E} on a fine $(\delta,x)$ grid, \eqref{eq:F}, $b\ge0$ and $h\le1$ (a linear program), one obtains $B=0.8749571$ at $\rho=0$ (Liu's $B_{\mathrm{alg}}$, the value at $\kappa=3/4$, which validates the setup), $0.873436$ at $\rho=0.01$, $0.872095$ at $\rho=0.02$, and $167/192$ for every $\rho\ge0.0404$, including $\rho=\infty$ (the saturation value is specific to a uniform shift and $\kappa=3/4$). Beyond that point the constraints \eqref{eq:E} are slack and the value is the bound of Proposition~\ref{prop:barrier}. To reach $B=5/6$ \emph{both} inputs must improve: the witness counts by $\rho\ge0.227$ (needed whatever is done to the floor class, since with the floor constraint deleted entirely $\rho_{\min}=1/6+\max_\delta(\alpha-\delta)\delta P_{1/2}/(2J)=0.2271$ at $M=2/3$, $\ell=1/3$, $h=1$), and the floor-class exponent from $1$ to at most $61/75\approx0.813$ (a saving $U^{-14/75}$). Neither alone suffices: with the present witness counts, lowering the floor exponent even to $0$ leaves $B=0.8749571$, because the floor constraint is slack at Liu's optimum.
\end{remark}

\begin{remark}[where $R^*=2/3$ comes from]
At the critical class $(\delta_c,x=1/2)$ the proven count is a three-way balance of the inverse-polynomial moment \cite[Lemma~17.1]{Ope26} (capacity $(1-r)/2$), the plain fourth moment \cite[Lemma~18.1]{Ope26} (capacity $(1-2m)/(6\kappa)$, $\kappa=3/4$) and the amplification slope $\alpha=5/6$ \cite[Lemma~17.6]{Ope26}, balanced at the detector cutoff $t^*=1+\delta P_x/(2J)\approx1.124$. The sextic large sieve of \cite[Lemma~9.1]{Ope26} is not used in the $7/8$ stage. Linearised at the critical class, $\partial R^*/\partial(\text{inverse capacity scale})=-0.031$, $\partial R^*/\partial(\text{plain capacity scale})=-0.013$, $\partial R^*/\partial\alpha=+0.031$ (more capacity lowers $R^*$, a steeper amplification slope raises it): each input must change by an $O(1)$ relative amount to move $R^*$ by $0.04$. The moment parameter $\kappa$ is a very weak lever ($F(1/2)=0.874708$ in Liu's notation).
\end{remark}

\begin{remark}[the endpoint certificate]
Lemma~20.2 of \cite{Ope26} certifies $-E^*\ge49/440640\approx1.11\cdot10^{-4}$. The true minimum of $-E^*$ over $[1/50,5/6]\times[0,1/2]$ is $2.28\cdot10^{-4}$ at $(\delta,x)=(0.3866,1/2)$; the certificate is valid but not sharp, the slack coming from the bound $J\le5/2$ in its proof ($J\approx1.22$ at the minimiser; the other bound used there, $v\le17(3+5y)$, is tight at $y=0$).
\end{remark}

\section{The problem this isolates, and why it is not a mean-value problem}

Within the comparison system, nothing below $167/192=0.8698\ldots$ is available by any choice of parameters or any improvement of the zero-witness estimates; the October 2026 refinements \cite{Liu26,XS26a,Liu26b} all remain above it, and \cite{Liu26} treats the lowest class only by its count (\cite[Lemma~19.2(4)]{Liu26}; we have checked this in \cite{Liu26} only). What would remove the barrier is a non-trivial treatment of the floor class. It is important to state correctly what that would have to be.

\emph{Absolute values cannot work.} By \cite[Lemma~8.1]{Ope26}, which is stated for every retained row (the floor bin included; cf.\ \cite[\S10.5]{Ope26}), for every $\psi$ in the finite twist family $X_u$ one has $|L_{\mathrm{orig}}(s,\psi)|+|L_{\mathrm{orig}}(s,\psi)^{-1}|\ll U^{\varepsilon}$ for $\Re s\ge a+6e$, $|\Im s|\le(3i+2)T_1$, and $|L_{\mathrm{orig}}(1-a-6e+it,\psi)|\ll U^{a-1/2+12e+\varepsilon}(3+T_1)^C$; here $L_{\mathrm{orig}}$ differs from $L^S$ by a bounded two-sided factor and the height factor is absorbed by the choice of $\tau$ in \cite[\S11]{Ope26}. The retained lines $\Re s=a_0+16e$, $\Re w=1-a_0-6e$ lie in this region. Consequently, on the floor class the scalar quotient $L^S(w,\chi_\bullet(u))/L^S(s,\eta\chi_\bullet(u))$ has absolute value $U^{\delta_0/2+o(1)}$ for all but $O(U^{1-\theta/2})$ rows (those whose powerful part exceeds $U^{\theta}$, for which the conductor $Q_\psi\asymp q_{\mathrm{rad}\,u}$ is smaller than $q_u$); the two-sided bound comes from Lemma~8.1 together with the functional equation in its proof. The local correction $H_{\eta,u}$ satisfies $H_{\eta,u}\ll U^{\varepsilon}$ \cite[Lemma~7.1]{Ope26}, and no lower bound for it is asserted in \cite{Ope26} (Definition~10.1 there explicitly assumes none); the argument here concerns the scalar quotient. By \cite[Prop.~9.2, (9.16)]{Ope26} (and the sharper no-slot count of \cite[Prop.~19.2]{Ope26} for $\delta\le5/6$) every bin with $a>a_0$ has $\ll U^{R(\delta)+\varepsilon}(1+T_1)^A$ rows with $R(\delta)\le1-1/75$, there are $O(\log U)$ bins, and a dyad has $\asymp U$ retained rows; hence the floor class contains all but $O(U^{1-1/75+\varepsilon})$ of the rows. Therefore any bound that multiplies per-row absolute values of the scalar quotient by a count --- a negative moment of $L$, a large-sieve mean square of the inverse polynomial (whose diagonal is exactly the trivial count), or a zero-density estimate combined with the $U^{o(1)}$ bound for $\log L$ in a zero-free rectangle --- reproduces the trivial floor count up to $U^{o(1)}$. In particular the first negative moment $\sum_{u\ \mathrm{floor}}|L^S(s,\eta\chi_\bullet(u))|^{-1}$ is $\gg U^{1-\varepsilon}$ on the retained line, since Lemma~8.1 gives $|L^S(s,\eta\chi_\bullet(u))|\ll U^{\varepsilon}$ for each floor row. The literature on negative moments (Gonek; Bui--Florea; Gao; Gao--Zhao) gives asymptotics or bounds of the order of the family size, conditional on RH or GRH in the number-field case, and offers nothing here.

\emph{What is actually required.} The architecture bounds each bin by the triangle inequality applied to the row sum, with the reciprocal taken out pointwise \cite[proof of Lemma~10.4]{Ope26}; the hypothesis (10.13) there excludes the reciprocal. To lower the floor exponent one must instead exhibit cancellation in the \emph{signed} sum over floor rows of one bin,
\[
S_{\mathrm{floor}}(s,w,z)=\sum_{\substack{u\ \text{floor},\ q_u\asymp U}}\xi(u)\,q_u^{-z}\,\frac{L^S(w,\chi_\bullet(u))\,H_{\eta,u,Z}(s,w,z)}{L^S(s,\eta\chi_\bullet(u))}\ \ll\ U^{1-\rho+\delta_0/2+\varepsilon}Z^{\ell(z_0-1/2)+g}
\]
(the scalar $\zeta_F^S(6z)$ omitted), uniformly on the retained lines $\Re s=a_0+16e$, $\Re w=1-a_0-6e$, $\Re z=17/50$ and heights $\le T_1/2$, $T_1=Z^{\tau}\le U^{1/100}$, and replace (10.13) by a hypothesis on the full signed row sum including the reciprocal. Any $\rho>0$ lowers the bound of Proposition~\ref{prop:barrier} once the witness counts have improved by $\ge0.0404$ as in Remark~\ref{prop:barrier}; $\rho\ge14/75$ is what $B=5/6$ requires, together with the witness improvement $\ge0.227$. Two obstructions are visible. The floor class is defined by zeros, not arithmetically, so the only route to a signed sum over it is ``all rows minus detected rows''; but a signed sum over \emph{all} rows of $1/L^S(s,\eta\chi_\bullet(u))$ on $\Re s=a_0$ presupposes that every member of the family is zero-free there, a hypothesis far stronger than the conclusion. Replacing $1/L$ by a truncated inverse polynomial of length $D$, which is entire and valid for all rows, leaves (by Lemma~4.9 of \cite{Ope26} on the zero-free box of a floor row and a smoothed Perron formula shifted to $\Re(s+v)=a_0+6e$) a per-row error $U^{\varepsilon}D^{-10e}$, so that absolute summation over $\asymp U$ rows requires $D\ge U^{\rho/(10e)}$, i.e.\ $D\ge U^{18}$ for $e<10^{-3}$ and $\rho=14/75$ ($e$ cannot be enlarged much, since $(16-6l_y)e$ is charged in (10.16) against a total saving of order $1/24$); this is outside the range of \cite[Lemma~17.1]{Ope26} and useless in \cite[Lemmas~17.6, 9.1]{Ope26}, which give no saving at such lengths. In the expansion of the all-rows sum into Dirichlet series the diagonal terms, those with $mn$ a sixth power, carry the factor $\sum_u\xi(u)q_u^{-z}\mathbf 1_{(u,mnS)=1}(\cdots)$, which has no pole since $\xi$ is a nonprincipal primitive character modulo $b^*$ \cite[\S6.1]{Ope26}; so no main term of size $U$ arises in the all-rows sum. This is a heuristic about the all-rows sum and says nothing by itself about $S_{\mathrm{floor}}$. We therefore regard the barrier as intrinsic to the comparison system of \S1 with the floor class counted trivially; removing it requires the signed cancellation above, for which we found no route within the present architecture.

\begin{thebibliography}{9}
\bibitem{Tom26} tomoto0, Barrier analysis (\emph{障壁分析}), \texttt{reports/barrier\_analysis.md} in \texttt{github.com/tomoto0/quasi-riemann-hypothesis-7-8-verification}, commit \texttt{b83a73c}, 2026-10-07.
\bibitem{Tracker} GTimG, QuasiRiemannTracker, milestone ``13/15'', \texttt{github.com/GTimG/QuasiRiemannTracker}, \texttt{src/bound-milestones.ts}.
\bibitem{Liu26} B.~Liu, Improved zero-free half-planes for the quasi-Riemann hypothesis, arXiv:2610.12234v2 (2026), Lean-formalized.
\bibitem{Liu26b} B.~Liu, Zero-free half-planes for Hecke $L$-functions over number fields, arXiv:2610.26835 (2026).
\bibitem{XS26a} Y.~Xu, H.~Sang, Amplified moments and a refinement of the quasi-Riemann hypothesis, arXiv:2610.15291 (2026).
\bibitem{XS26b} Y.~Xu, H.~Sang, A matching lower bound for the sextic large sieve and structured prime moments, arXiv:2610.17675 (2026).
\bibitem{Ope26} OpenAI, The quasi-Riemann hypothesis: a zero-free half-plane $\Re s>7/8$, manuscript (2026-09-30), \texttt{github.com/openai/math}, family 003, Lean-verified.
\end{thebibliography}
\end{document}
